
% SGA 3，报告 VII_B；英文重建本第一轮的中文译本。
% 范围： 节 3.3.3 全文 依据文献局部页 72--73,
% 止于下节之前： 节 4.
% 合并阅读本页： 585--586. 重版印刷页： 574--575.
% 范围内原版边码： 539--540.
% 范围内编者注： 132--133.
% 依据： Polo--Gille Exp7B-13oct24.pdf,
% SHA-256 D303690E0195654BBE7C3DAA5D89EA23E24CD709A3575217234AC19AE8836694.
% 参照了 Jacob Reinhold 的 CC BY 4.0 Markdown，作为对照来源。
% 原英文 Loop 1 阶段说明：纤维积图为第一轮直接从 PDF 截取的原图。（当前中文工程已采用原生图。）
\subsubsection*{3.3.3}
\phantomsection
\label{sga3:VIIB:3.3.3}
继续设 \(k\) 是特征 \(0\) 的域。令 \(\overline k\) 为 \(k\) 的代数闭包，\(\Gamma\) 为 \(\overline k\) 在 \(k\) 上的拓扑 Galois 群。
对每个形式 \(k\)-群 \(G\)，用 \(\mathbf{Aut}_k(G)\) 表示将每个有限维交换 \(k\)-代数 \(C\) 映为形式 \(C\)-群
\[
  G\widehat\otimes_k C.
\]
的自同构群的 \(k\)-群函子。\footnote{\label{sga3:VIIB:nde:132}原文说：“这个 \(k\)-函子显然是左正合的\textup{（\(H\) 在 \(k\) 上拓扑平坦！）}。”以下细节由编者补充。}
由于 \(k\) 是域，\(G\) 在 \(k\) 上拓扑平坦。等价地，其仿射代数
\[
  A=\mathcal A(G)
\]
在 \(k\) 上拓扑平坦，或者其协变双代数
\[
  H=\mathbf H(G)
\]
在 \(k\) 上平坦。因此函子 \(\mathbf{Aut}_k(G)\) 是形式 \(k\)-群。
事实上，\(\mathbf{Aut}_k(G)\) 是下图的纤维积\textup{（参见\ \SGARefLink{sga3:i:remark:1-7-3}{报告 I，1.7.3}）}：
\begin{center}
  \phantomsection
  \label{sga3:VIIB:3.3.3:automorphism-fiber-product-diagram}
  % SGA3-VII-DIAG-120
  % SGA3-VII-DIAG-120
% Authority: Exp7B-13oct24.pdf, local p. 72.
% Crop: figures/Exp7B_localp072_automorphism_fiber_product_diagram.png
% Locator: tex/components/63_expose_VIIB_section333_gamma_lie_algebras_en.tex, Loop 1 line 49.
\begin{tikzcd}[row sep=3.2em,column sep=3.8em]
{} &
\underline{\operatorname{End}}_k(G)
  \times\underline{\operatorname{End}}_k(G) \arrow[d] \\
\operatorname{Spf}(k) \arrow[r] &
\underline{\operatorname{End}}_k(G)
  \times\underline{\operatorname{End}}_k(G)
\end{tikzcd}

\end{center}
竖直箭头将元素映为
\[
  (\phi,\psi)\longmapsto(\phi\circ\psi,\psi\circ\phi),
\]
而水平箭头是单位截面。因此，只须证明 \(k\)-函子 \(\mathbf{End}_k(G)\) 可由 \(\operatorname{Vaf}/k\) 的对象表示；由 \SGARefLink{sga3:VIIB:1.1}{1.1} 和 \SGARefLink{sga3:VIIB:0.4.2}{0.4.2}，这归结为证明它是左正合的，或者等价地，它与 \(k\)-代数的纤维积交换。
给定 \(\mathbf{End}_k(G)(C)\) 的一个元素，等价于给定一个同时保持余代数结构的 \(k\)-代数态射
\[
  \phi:H\longrightarrow H\otimes_k C
\]
也就是说，通常的相容图对该态射交换。由于 \(H\) 在 \(k\) 上平坦，函子
\[
  H\otimes_k -
\]
与 \(k\)-代数的纤维积交换。因此 \(\mathbf{End}_k(G)\) 是左正合的，并且可视为形式 \(k\)-群。
若 \(L\) 是 \(k\)-Lie 代数，\(G=\mathscr G(L)\)，则\SGARefLink{sga3:VIIB:3.3:cartier-theorem}{定理 3.3} 表明，\(\mathbf{Aut}_k(G)\) 同构于 \(k\)-群函子 \(\mathbf{Aut}_k(L)\)，后者将每个有限维 \(k\)-代数 \(C\) 映为 \(C\)-Lie 代数
\[
  C\otimes_k L.
\]
的自同构群。
现在设 \(G\) 是任意形式 \(k\)-群。\SGARefLink{sga3:VIIB:2.5.2:proposition}{2.5.2} 已表明，\(G\) 典范分解为平展形式群 \(G_e\) 与无穷小形式群 \(G_0\) 的半直积。这个半直积由以下数据确定：
\[
  L=\operatorname{Lie}(G)=\operatorname{Lie}(G_0),
\]
、\(\Gamma\)-群
\[
  M=G_e(\overline k)=G(\overline k),
\]
以及形式 \(k\)-群态射
\[
  \Phi:
  G_e\longrightarrow
  \mathbf{Aut}_k(G_0)
  \simeq
  \mathbf{Aut}_k(L).
\]
这样的态射将 \(G_e\) 映入 \(\mathbf{Aut}_k(L)\) 的平展部分 \(\mathbf{Aut}_k(L)_e\)\textup{（参见\ \SGARefLink{sga3:VIIB:2.5.1:corollary}{2.5.1}）}。因此，它由 \(\Gamma\)-群态射
\[
  \phi=\Phi(\overline k):
  M\longrightarrow
  \bigl(\mathbf{Aut}_kL\bigr)(\overline k)
  =
  \operatorname{Aut}_{\overline k}
    \bigl(L\otimes_k\overline k\bigr).
\]
确定。
令 \(\Gamma\) 如下作用于 \(L\otimes_k\overline k\)
\[
  \gamma(\ell\otimes t)=\ell\otimes\gamma(t).
\]
则 \(\phi\) 使 \(M\) 通过 \(\overline k\)-Lie 代数自同构作用于 \(L\otimes_k\overline k\)，使得
\[
  \phi\bigl(\gamma(m)\bigr)
  =
  \gamma\circ\phi(m)\circ\gamma^{-1}.
\]
等价地，
\[
  \gamma(m)(\ell\otimes t)
  =
  \gamma\!\left(
    m\bigl(\ell\otimes\gamma^{-1}(t)\bigr)
  \right)
\]
对每个 \(m\in M\)、\(\gamma\in\Gamma\) 和 \(\ell\otimes t\in L\otimes_k\overline k\) 成立。我们将最后这个条件表述为：\(M\) 在 \(L\otimes_k\overline k\) 上的作用\emph{与 \(\Gamma\) 相容}。
可将这一情形概括成一个“抽象”的陈述。所谓 \emph{\(k\) 上的 \(\Gamma\)-Lie 代数}，是三元组
\[
  (L,M,\phi)
\]
由 \(k\)-Lie 代数 \(L\)、\(\Gamma\)-群 \(M\)，以及 \(M\) 在 \(L\otimes_k\overline k\) 上与 \(\Gamma\) 在 \(M\) 和 \(L\otimes_k\overline k\) 上的作用相容的作用 \(\phi\) 组成。
若 \((L,M,\phi)\) 和 \((L',M',\phi')\) 是两个这样的 \(\Gamma\)-Lie 代数，则从前者到后者的态射是一个对 \((f,\theta)\)，其中
\[
  f:L\longrightarrow L'
\]
是 \(k\)-Lie 代数态射，而
\[
  \theta:M\longrightarrow M'
\]
是 \(\Gamma\)-群态射，使得
\[
  (f\otimes1)(m\cdot x)
  =
  \theta(m)\cdot(f\otimes1)(x)
\]
对每个 \(x\in L\otimes_k\overline k\) 和 \(m\in M\) 成立。因此，我们证明了：
\medskip
\noindent\textbf{定理。}
\phantomsection
\label{sga3:VIIB:3.3.3:gamma-lie-equivalence}
{\itshape 设 \(k\) 是特征 \(0\) 的域。将形式 \(k\)-群 \(G\) 映为三元组
\[
  \bigl(
    \operatorname{Lie}(G),
    G(\overline k),
    \Phi(\overline k)
  \bigr)
\]
的函子，给出形式 \(k\)-群范畴与 \(\Gamma\)-Lie 代数范畴之间的等价。}\footnote{%
\label{sga3:VIIB:nde:133}特别地，当 \(k=\overline k\) 时，就重新得到 \SGARefLink{sga3:VIIB:2.9.2}{2.9.2} 中提到的 Cartier--Gabriel--Kostant 定理。}
